% part: intuitionistic-logic
% chapter: propositions-as-types
% section: normalization

\documentclass[../../../include/open-logic-section]{subfiles}

\begin{document}

\olfileid{int}{pty}{nor}

\olsection{正规化}

本节我们将证明，通过某种归约顺序，任何推导都能归约为正规推导，这被称为\emph{正规化性质}。
我们将利用命题即类型对应：证明每个证明项都能归约为正规形式；自然演绎推导的正规化则可随之得出。

首先我们定义一些度量项复杂度的函数。
公式的\emph{长度}~$\len{!A}$ 定义为
\begin{align*}
  \len{p} &= 0\\
  \len{!A \land !B} &= \len{!A} + \len{!B} + 1\\
  \len{!A \lor !B} &= \len{!A} + \len{!B} + 1 \\
  \len{!A \lif !B} &= \len{!A} + \len{!B} + 1.
\end{align*}
可约式~$M$ 的复杂度由其\emph{切秩}~$\cutrank{M}$ 度量：
\begin{align*}
  \cutrank{(\lambd[\typeof{x}{!A}][\typeof{N}{!B}])Q} &=
  \len{!A} + \len{!B} + 1 \\
  \cutrank{\ande{i}{\andi{\typeof{M}{!A}}{\typeof{N}{!B}}}} &=
  \len{!A} + \len{!B} + 1 \\
  \cutrank{\ore{\ori{i}{!A_{i}}{\typeof{M}{!A_i}}}
    {\typeof{x_1}{!A_1}}{\typeof{N_1}{!C}}
    {\typeof{x_2}{!A_2}}{\typeof{N_2}{!C}}} &=
  \len{!A} + \len{!B} + 1
\end{align*}
证明项的复杂度由其内部最复杂的可约式度量，若为正规形式则为 $0$：
\[
  \maxrank{M} = \max \{\cutrank{N} | N \text{ 是 } M \text{ 的子项且是可约式}\}
\]

\begin{lem}\ollabel{lem:subst}
  若 $\Subst{M}{\typeof{N}{!A}}{\typeof{x}{!A}}$ 是一个可约式且 $M
  \not\ident x$，则下列情形之一成立：
  \begin{enumerate}
  \item $M$ 本身是一个可约式，或
  \item $M$ 形如 $\ande{i}{x}$，且 $N$ 形如
    $\andi{P_1}{P_2}$
  \item $M$ 形如 $\ore{i}{x_1}{P_1}{x_2}{P_2}$，且 $N$ 形如
    $\ori{i}{}{Q}$
  \item $M$ 形如 $x Q$，且 $N$ 形如
    $\lambd[x][P]$
  \end{enumerate}
  在第一种情形中，$\cutrank{\Subst{M}{N}{x}} = \cutrank{M}$；在其他情形中，$\cutrank{\Subst{M}{N}{x}} = \len{!A})$。
\end{lem}
\begin{proof}
  对~$M$ 进行归纳证明。
  \begin{enumerate}
  \item 若 $M$ 是单个变元 $y$ 且 $y \not\ident x$，则
    $\Subst{y}{N}{x}$ 就是 $y$，因此不是可约式。
  \item 若 $M$ 形如 $\andi{N_1}{N_2}$，或 $\lambd[x][N]$，或
    $\ori{i}{!A}{N}$，那么 $\Subst{M}{\typeof{N}{!A}}{\typeof{x}{!A}}$ 也
    具有相应的形式，因此不是可约式。
  \item 若 $M$ 形如 $\ande{i}{P}$，我们分两种情形考虑。
    \begin{enumerate}
      \item 若 $P$ 形如 $\andi{P_1}{P_2}$，那么 $M \ident
        \ande{i}{\andi{P_1}{P_2}}$ 是一个可约式，且显然
        \[
        \Subst{M}{N}{x} \ident
        \ande{i}{\andi{\Subst{P_1}{N}{x}}{\Subst{P_2}{N}{x}}}
        \]
        也是一个可约式。二者切秩相等。
      \item 若 $P$ 是单个变元，则它必须是 $x$ 才能使代换后的结果为可约式，且 $N$ 必须形如
        $\andi{P_1}{P_2}$。现在考虑
        \[
        \Subst{M}{N}{x} \ident \Subst{\ande{i}{x}}{\andi{P_1}{P_2}}{x},
        \]
        它是 $\ande{i}{\andi{P_1}{P_2}}$。它的切秩等于
        $\cutrank{x}$，即 $\len{!A}$。
    \end{enumerate}
  \end{enumerate}
  $\ore{N}{x_1}{N_1}{x_2}{N_2}$ 和 $P Q$ 的情形类似。
\end{proof}

\begin{lem}
  若 $M$ 归约为 $M'$，且对 $M$ 的所有真可约式子项 $N$ 都有 $\cutrank{M} > \cutrank{N}$，则 $\cutrank{M} >
  \maxrank{M'}$。
\end{lem}

\begin{proof}
  分情形证明。
  \begin{enumerate}
  \item 若 $M$ 形如 $\ande{i}{\andi{M_1}{M_2}}$，则 $M'$ 是
    $M_i$；由于 $M_i$ 的任何子项也是 $M$ 的真子项，故断言成立。
  \item 若 $M$ 形如
    $(\lambd[\typeof{x}{!A}][N])\typeof{Q}{!A}$，则 $M'$ 是
    $\Subst{N}{\typeof{Q}{!A}}{\typeof{x}{!A}}$。 考虑 $M'$ 中的一个可约式。
      要么 $N$ 中存在与之对应且切秩相等的可约式，根据假设该秩小于 $\cutrank{M}$；要么其切秩等于 $\len{!A}$，根据定义它小于
    $\cutrank{(\lambd[x^{!A}][N])Q}$。
  \item 若 $M$ 形如
    \[
    \ore{\ori{i}{}{\typeof{N}{!A_i}}}
        {\typeof{x_1}{!A_1}}{\typeof{N_1}{!C}}
        {\typeof{x_2}{!A_2}}{\typeof{N_2}{!C}},
    \]
    则 $M' \ident \Subst{N_i}{N}{\typeof{x_i}{!A_i}}$。 考虑
    $M'$ 中的一个可约式。要么 $N_i$ 中存在与之对应且切秩相等的可约式，根据假设该秩小于 $\cutrank{M}$；要么其切秩等于 $\len{!A_i}$，根据定义它小于
    $\cutrank{\ore{\ori{i}{}{\typeof{N}{!A_i}}}
      {\typeof{x_1}{!A_1}}{\typeof{N_1}{!C}}
      {\typeof{x_2}{!A_2}}{\typeof{N_2}{!C}}}$。
  \end{enumerate}
\end{proof}

\begin{thm}
  所有证明项均可归约为正规形式；所有推导均可归约为正规推导。
\end{thm}

\begin{proof}
  第二个结论由第一个结论得出。我们通过以 $m=\maxrank{M}$（其中 $M$ 是一个证明项）为参数的完全归纳法来证明第一个结论。
  \begin{enumerate}

  \item 若 $m=0$，则 $M$ 已为正规形式。
  \item
    否则，我们对 $M$ 中切秩恰为 $m$ 的可约式的数目 $n$ 进行归纳。
    \begin{enumerate}
    \item 若 $n=1$，选取任意一个满足“对 $M$ 的任何同为可约式的真子项 $P$（当然如此），都有 $m = \cutrank{N} >
      \cutrank{P}$”的可约式 $N$。这样的可约式必然存在，因为任何项都只有有限多个子项。

      记 $N'$ 为 $N$ 的约化式。根据引理，有 $\maxrank{N'} < \maxrank{N}$，因此，$n$（即切秩为 $m$ 的可约式数目）减少了。所以 $m$ 减少了（减少 $1$ 或更多），从而可以对 $m$ 应用归纳假设。
    \item 对于归纳步骤，假设 $n>1$。过程类似，区别在于 $n$ 只是减少到一个正数，因而 $m$ 并未改变。我们直接对 $n$ 应用归纳假设即可。
    \end{enumerate}
\end{enumerate}
\end{proof}

类型化 $\lambd$-项可正规化这一事实，意味着每个这样的项都\emph{有一个正规形式}。如果我们把 $\lambd$-项 $\beta$-归约到其正规形式看作执行一次计算，并把正规形式看作计算的值，那么这就意味着类型化 $\lambd$-演算中的每次计算最终都会停机并产生一个值。此外，类型保持性保证了该值具有正确的类型。例如，如果 $N$ 具有类型 $!A \lif !B$（即它是一个函数，旨在处理类型为~$!A$ 的自变元并产生类型为~$!B$ 的值），并且我们将其应用于一个类型为~$!A$ 的项~$M$（作为输入），那么 $(NM)$ 会归约为一个项 $N'$（作为输出），并且能保证该输出的类型是~$!B$。

事实上，类型化 $\lambd$-演算中的项不仅仅对将 $\beta$-归约应用于子项的某一种方式会归约到正规形式，而且对\emph{任何}将 $\beta$-归约应用于子项的方式，最终都会终止于正规形式。这一性质被称为\emph{强正规化}。它不仅意味着存在一种方式能保证计算最终停机并产生一个值，而且意味着\emph{任何}执行计算的方式最终都会产生一个值。

我们如何知道不同的计算执行方式会产生\emph{相同的}值呢？到目前为止，我们只知道（由类型保持性可知）它产生的所有值都具有相同的类型。类型化 $\lambd$-演算还有另一个重要性质：它是\emph{合流的}，或者说具有 \emph{邱奇-罗瑟} 性质。如果 $N \red N_1$ 且 $N \red N_2$，那么存在某个项~$N'$ 使得 $N_1 \red N'$ 且 $N_2 \red N'$。回想一下，$N \red N'$ 表示：$N'$ 是通过零次或多次 $\redone$-归约步骤得到的。如果 $N_1$ 是正规形式，那么使得 $N_1 \red N'$ 的唯一项 $N'$ 就是 $N_1$ 自身（通过零次 $\redone$-归约步骤）。由此可知，如果 $N_1$ 和 $N_2$ 都是正规形式，那么 $N_1 = N' = N_2$。换句话说，既然 $\red$ 是强正规化的，那么归约一个项的任何方式最终都会产生一个正规形式。既然 $\red$ 是 邱奇-罗瑟 的，那么任意两个正规形式都是相等的。因此，类型化 $\lambd$-演算中的计算总是会终止，并且无论计算如何进行，得到的值（正规形式）都是相同的。

证明一个关系是合流的通常很困难。但我们可以容易地确立以下更弱的条件：

\begin{prop}[弱 Church-Rosser 性质]
如果 $N \redone N_1$ 且 $N \redone N_2$，那么存在项 $N'$ 使得 $N_1 \red N'$ 且 $N_2 \red N'$。
\end{prop}

\begin{proof}
项 $N_1$ 和 $N_2$ 分别是将 $N$ 中的可约式 $M_1$ 或 $M_2$ 替换为其约化式而得到的。可约式 $M_1$ 和 $M_2$ 是 $N$ 的子项，因此不会部分重叠。这意味着要么一个项是另一个项的子项，要么它们出现在 $N$ 中不同的、互不重叠的位置。考虑后一种情况：$N$ 形如 $N(M_1,M_2)$。将 $M_1$ 替换为其约化式 $M_1'$ 得到 $N_1$，即 $N(M_1',M_2)$；将 $M_2$ 替换为其约化式 $M_2'$ 得到 $N_2$，即 $N(M_1,M_2')$。两者都归约到 $N(M_1',M_2')$。对于另一种情况，假设 $M_2$ 是 $M_1$ 的一个子项（另一种情况对称），并根据 $M_1$ 是哪一类可约式来区分情况。例如，如果 $M_1 \ident
  \proj{1}{\pair{O_1}{O_2}}$，它归约为 $O_1$。如果 $M_2$ 是 $O_1$ 的一个子项，那么归约 $M_2$ 会得到 $O_1'$。因此，先归约 $M_1$ 再归约 $M_2$ 会得到 $O_1'$。但是先归约 $M_2$ 会得到 $\proj{1}{\pair{O_1'}{O_2}}$，而这也归约为 $O_2$。
\end{proof}

\begin{lem}[Newman 引理]
如果 $\red$ 是强正规化的并且具有弱 邱奇-罗瑟 性质，那么它是合流的。
\end{lem}

\begin{proof}
由于 $\red$ 是强正规化的，任何归约序列都会产生一个正规形式。正规形式唯一，当且仅当 $\red$ 是合流的。因此，假设某个项 $N$ 有两个不同的正规形式 $N'$ 和 $N''$，分别来自归约序列
\begin{align*}
  N & = N_1' \redone N_2' \redone \dots \redone N_k' = N' \text{ 和}\\
  N & = N_1'' \redone N_2'' \redone \dots \redone N_l'' = N''
  \end{align*}
（归约序列的长度必须 $>0$，否则 $N$ 本身就已经是正规形式了）。

如果 $N_1' = N_1''$，那么它（即 $N_1'$）就有两个不同的正规形式（$N'$ 和 $N''$）。如果 $N_1' \neq N_1''$，那么由弱 邱奇-罗瑟 性质可知，存在某个 $N'''$ 使得 $N_1' \red N'''$ 且 $N_1'' \red N'''$。由于 $N' \neq N''$，要么 $N''' \neq N'$，要么 $N''' \neq N''$。因此，要么 $N_1'$ 要么 $N_1''$ 有两个不同的正规形式。我们已经证明了：如果 $N$ 有两个不同的正规形式，那么就存在一个 $N_1$ 使得 $N \redone N_1$ 且 $N_1$ 有两个不同的正规形式。显然我们可以继续这一过程并得到 $N \redone N_1
  \redone N_2 \redone \dots$，但这是不可能的，因为 $\redone$ 是强正规化的。
\end{proof}


\end{document}